% Parameter pemetaan skor 0-2 perbandingan kematangan MLOps
\begin{table}[!htb]
\centering
\caption{Parameter Pemetaan Skor Perbandingan Kematangan MLOps}
\label{tab:rubrik_maturity}
\footnotesize
\setlength{\tabcolsep}{3pt}
\renewcommand{\arraystretch}{1.2}
\begin{tabular}{|>{\centering\arraybackslash}m{1.2cm}|>{\raggedright\arraybackslash}m{12.1cm}|}
\hline
\centering\arraybackslash\textbf{Nilai} & \centering\arraybackslash\textbf{Parameter Pemetaan dari Deskripsi \textit{Level} Terbit} \\
\hline
0 & Praktik pada kolom tidak disebut sama sekali pada deskripsi terbit \textit{level} tersebut, atau disebut secara eksplisit belum ada. \\
\hline
1 & Praktik disebut hadir sebagian, berjalan semi-otomatis, atau berstatus opsional pada deskripsi terbit \textit{level} tersebut. \\
\hline
2 & Praktik disebut berjalan otomatis penuh atau menjadi komponen wajib pada deskripsi terbit \textit{level} tersebut. \\
\hline
\end{tabular}
\end{table}
